\chapter{PHP et MySQL : les bases de données}

\begin{synthese}[Au programme]
Stocker durablement des données dans une base relationnelle MySQL, écrire les
requêtes SQL essentielles (\code{CREATE}, \code{INSERT}, \code{SELECT},
\code{UPDATE}, \code{DELETE}), puis interfacer cette base depuis PHP avec
l'extension \code{mysqli} : se connecter, exécuter une requête, lire et afficher
les résultats. C'est le cœur de toute application web de gestion.
\end{synthese}

\section{La base de données relationnelle}

Jusqu'ici, dès que le script PHP se termine, les données disparaissent. Pour
\textbf{conserver} les notes d'un élève, la liste des produits d'une boutique de
Douala ou les inscriptions à un concours, on a besoin d'une \textbf{base de
données}. Le logiciel qui la gère s'appelle un \textbf{SGBD} (Système de Gestion
de Base de Données) ; ici, ce sera \textbf{MySQL}.

\begin{definition}
Une \textbf{base de données relationnelle} organise l'information dans des
\textbf{tables}. Une table ressemble à un tableau à deux dimensions :
\begin{itemize}
  \item une \textbf{ligne} (ou \emph{enregistrement}) représente \textbf{un}
  élément concret (un élève, un produit\dots) ;
  \item une \textbf{colonne} (ou \emph{champ}) représente une \textbf{propriété}
  commune à tous les éléments (le nom, la classe, la moyenne\dots).
\end{itemize}
\end{definition}

\begin{exemple}
La table \code{eleves} du Lycée de Maroua :

\begin{center}\small
\begin{tabular}{@{}llll@{}}
\toprule
\textbf{id} & \textbf{nom} & \textbf{classe} & \textbf{moyenne}\\
\midrule
1 & Awa Bello   & Tle TI & 14.50\\
2 & Jean Mballa & Tle TI & 9.75\\
3 & Sandrine Eto & Tle TI & 16.00\\
\bottomrule
\end{tabular}
\end{center}
\fig{Trois enregistrements (lignes) et quatre champs (colonnes).}
\end{exemple}

\begin{definition}
La \textbf{clé primaire} (\emph{primary key}) est un champ dont la valeur est
\textbf{unique} pour chaque ligne : elle identifie un enregistrement sans
ambiguïté. On utilise très souvent un champ \code{id} entier qui
\textbf{s'incrémente automatiquement} (\code{AUTO\_INCREMENT}).
\end{definition}

\begin{remarque}
Deux élèves peuvent porter le même nom, mais jamais le même \code{id}. La clé
primaire sert justement à distinguer puis à retrouver précisément une ligne, par
exemple pour la modifier ou la supprimer.
\end{remarque}

\section{Le langage SQL}

On dialogue avec MySQL grâce au langage \textbf{SQL} (\emph{Structured Query
Language}). Une instruction SQL s'appelle une \textbf{requête}. Par convention,
on écrit les mots-clés en \textbf{MAJUSCULES} et on termine chaque requête par un
point-virgule \code{;}.

\subsection{Créer une table : \texttt{CREATE TABLE}}

\begin{syntaxe}
\code{CREATE TABLE} \textit{nom\_table} \code{(} \textit{champ type contraintes}, \dots \code{);}
\end{syntaxe}

\begin{codesql}
CREATE TABLE eleves (
  id       INT AUTO_INCREMENT PRIMARY KEY,  -- clé primaire auto-incrémentée
  nom      VARCHAR(60) NOT NULL,            -- texte court, obligatoire
  classe   VARCHAR(20) NOT NULL,
  moyenne  DECIMAL(4,2) DEFAULT 0           -- nombre : 2 décimales (ex. 14.50)
);
\end{codesql}

\begin{remarque}
Les principaux types : \code{INT} (entier), \code{VARCHAR(n)} (chaîne d'au plus
\code{n} caractères), \code{DECIMAL(t,d)} (nombre décimal), \code{TEXT} (texte
long), \code{DATE} (une date). \code{NOT NULL} interdit une valeur vide.
\end{remarque}

\subsection{Insérer des données : \texttt{INSERT INTO}}

\begin{syntaxe}
\code{INSERT INTO} \textit{table} \code{(}\textit{champs}\code{) VALUES (}\textit{valeurs}\code{);}
\end{syntaxe}

\begin{codesql}
INSERT INTO eleves (nom, classe, moyenne)
VALUES ('Awa Bello', 'Tle TI', 14.50);

-- On ne précise pas id : MySQL le génère tout seul (1, 2, 3, ...)
INSERT INTO eleves (nom, classe, moyenne)
VALUES ('Jean Mballa', 'Tle TI', 9.75);
\end{codesql}

\begin{attention}
Les valeurs de type texte (\code{VARCHAR}, \code{TEXT}, \code{DATE}) se mettent
entre apostrophes simples \code{'\dots'}. Les nombres, eux, s'écrivent sans
apostrophes.
\end{attention}

\subsection{Lire des données : \texttt{SELECT}}

C'est la requête la plus utilisée : elle \textbf{interroge} la table et renvoie
les lignes demandées.

\begin{syntaxe}
\code{SELECT} \textit{champs} \code{FROM} \textit{table} \code{WHERE} \textit{condition} \code{ORDER BY} \textit{champ};
\end{syntaxe}

\begin{codesql}
SELECT * FROM eleves;                    -- toutes les colonnes, toutes les lignes

SELECT nom, moyenne FROM eleves;         -- seulement deux colonnes

SELECT * FROM eleves WHERE moyenne >= 10;        -- les admis

SELECT * FROM eleves WHERE classe = 'Tle TI';    -- filtre par classe

SELECT * FROM eleves ORDER BY moyenne DESC;      -- du plus fort au plus faible
\end{codesql}

\begin{remarque}
L'étoile \code{*} signifie « toutes les colonnes ». \code{ORDER BY \dots\ ASC}
trie en ordre croissant (par défaut), \code{DESC} en ordre décroissant. Le mot
\code{WHERE} introduit une \textbf{condition} de filtrage.
\end{remarque}

\subsection{Modifier des données : \texttt{UPDATE}}

\begin{syntaxe}
\code{UPDATE} \textit{table} \code{SET} \textit{champ = valeur} \code{WHERE} \textit{condition};
\end{syntaxe}

\begin{codesql}
UPDATE eleves SET moyenne = 15.25 WHERE id = 1;

UPDATE eleves SET classe = 'Tle TI2' WHERE classe = 'Tle TI';
\end{codesql}

\begin{attention}
\textbf{Ne jamais oublier le \code{WHERE} !} Sans condition,
\code{UPDATE eleves SET moyenne = 0;} mettrait la moyenne de \textbf{tous} les
élèves à zéro.
\end{attention}

\subsection{Supprimer des données : \texttt{DELETE}}

\begin{syntaxe}
\code{DELETE FROM} \textit{table} \code{WHERE} \textit{condition};
\end{syntaxe}

\begin{codesql}
DELETE FROM eleves WHERE id = 2;          -- supprime l'élève n°2

DELETE FROM eleves WHERE moyenne < 5;     -- supprime plusieurs lignes
\end{codesql}

\begin{attention}
Comme pour \code{UPDATE}, sans \code{WHERE} la requête \code{DELETE FROM eleves;}
vide \textbf{toute} la table. On supprime presque toujours par clé primaire
(\code{WHERE id = \dots}).
\end{attention}

\section{phpMyAdmin}

\begin{definition}
\textbf{phpMyAdmin} est une application web qui permet d'administrer une base
MySQL à la souris, depuis le navigateur, sans écrire (au début) de code. Livré
avec WampServer et XAMPP, on l'ouvre à l'adresse
\code{http://localhost/phpmyadmin}.
\end{definition}

On y crée une base, on y ajoute des tables et des lignes via des formulaires, et
surtout on dispose d'un onglet \textbf{SQL} où l'on tape directement les requêtes
vues plus haut pour les tester. C'est l'outil idéal pour préparer sa base
\textbf{avant} de l'interfacer avec PHP.

\begin{methode}
\textbf{Préparer sa base avec phpMyAdmin.} (1)~Ouvrir
\code{http://localhost/phpmyadmin} ; (2)~cliquer sur \emph{Nouvelle base de
données}, la nommer \code{lycee} ; (3)~onglet \emph{SQL} : coller le
\code{CREATE TABLE eleves} puis \emph{Exécuter} ; (4)~ajouter quelques lignes
avec \code{INSERT INTO}.
\end{methode}

\section{Se connecter à MySQL depuis PHP : \texttt{mysqli}}

Pour piloter MySQL depuis un script, PHP fournit l'extension \textbf{mysqli}
(\emph{MySQL improved}). On l'utilise ici en \textbf{style procédural} : de
simples fonctions \code{mysqli\_xxx()}.

\begin{definition}
\code{mysqli\_connect(\$hote, \$utilisateur, \$mot\_de\_passe, \$base)} ouvre une
connexion au serveur MySQL et renvoie un \textbf{identifiant de connexion}.
\end{definition}

En local (WampServer/XAMPP), l'hôte est \code{localhost}, l'utilisateur est
\code{root} et le mot de passe est en général \textbf{vide}.

\begin{codephp}
<?php
  // connexion.php — on se connecte une seule fois, puis on inclut ce fichier
  $hote = "localhost";
  $user = "root";
  $pass = "";          // mot de passe vide en local
  $base = "lycee";

  // Tentative de connexion
  $cnx = mysqli_connect($hote, $user, $pass, $base);

  // Gestion de l'erreur : si la connexion échoue, on arrête tout
  if (!$cnx) {
    die("Échec de la connexion : " . mysqli_connect_error());
  }

  // (Optionnel) accents corrects : on travaille en UTF-8
  mysqli_set_charset($cnx, "utf8");
?>
\end{codephp}
\fig{Le script \code{connexion.php}, réutilisé par toutes les pages.}

\begin{remarque}
La fonction \code{die("\dots")} (synonyme de \code{exit}) affiche un message puis
\textbf{stoppe} immédiatement le script. \code{mysqli\_connect\_error()} renvoie
la raison exacte de l'échec (mauvais mot de passe, base inexistante\dots).
\end{remarque}

\section{Exécuter une requête : \texttt{mysqli\_query}}

\begin{definition}
\code{mysqli\_query(\$cnx, \$requete)} envoie une requête SQL au serveur. Pour un
\code{SELECT}, elle renvoie un \textbf{jeu de résultats} ; pour
\code{INSERT}/\code{UPDATE}/\code{DELETE}, elle renvoie \code{true} en cas de
succès, \code{false} en cas d'erreur.
\end{definition}

\begin{codephp}
<?php
  require "connexion.php";   // on récupère $cnx

  $sql = "SELECT * FROM eleves";
  $res = mysqli_query($cnx, $sql);

  if (!$res) {
    die("Erreur de requête : " . mysqli_error($cnx));
  }
?>
\end{codephp}

\begin{remarque}
\code{require "connexion.php";} insère le fichier de connexion ; on récupère
ainsi la variable \code{\$cnx} sans réécrire le code à chaque page.
\code{mysqli\_error(\$cnx)} donne le détail d'une erreur de requête.
\end{remarque}

\section{Lire les résultats d'un \texttt{SELECT}}

Le jeu de résultats contient plusieurs lignes. On les parcourt \textbf{une par
une} avec une boucle \code{while} et la fonction \code{mysqli\_fetch\_assoc} :
elle renvoie la ligne suivante sous forme de \textbf{tableau associatif}
(\code{\$ligne["nom"]}, \code{\$ligne["moyenne"]}\dots), puis \code{false}
quand il n'y a plus de ligne.

\begin{syntaxe}
\code{while (\$ligne = mysqli\_fetch\_assoc(\$res)) \{ \dots\ \}}
\end{syntaxe}

\begin{codephp}
<?php
  require "connexion.php";
  $res = mysqli_query($cnx, "SELECT * FROM eleves ORDER BY nom");
?>
<table border="1">
  <tr><th>Nom</th><th>Classe</th><th>Moyenne</th></tr>
<?php
  // On parcourt chaque ligne du résultat
  while ($ligne = mysqli_fetch_assoc($res)) {
    echo "<tr>";
    echo "<td>" . $ligne["nom"]     . "</td>";
    echo "<td>" . $ligne["classe"]  . "</td>";
    echo "<td>" . $ligne["moyenne"] . "</td>";
    echo "</tr>";
  }
?>
</table>
\end{codephp}
\fig{Affichage de la table \code{eleves} dans un tableau HTML.}

\subsection{Les autres fonctions de lecture}

\begin{itemize}
  \item \code{mysqli\_fetch\_assoc(\$res)} : ligne en tableau \textbf{associatif}
  (\code{\$ligne["nom"]}) — la plus lisible ;
  \item \code{mysqli\_fetch\_row(\$res)} : ligne en tableau \textbf{numéroté}
  (\code{\$ligne[0]}, \code{\$ligne[1]}\dots) ;
  \item \code{mysqli\_fetch\_array(\$res)} : les deux à la fois (par index
  \textbf{et} par nom) ;
  \item \code{mysqli\_num\_rows(\$res)} : le \textbf{nombre} de lignes
  renvoyées.
\end{itemize}

\begin{codephp}
<?php
  require "connexion.php";
  $res = mysqli_query($cnx, "SELECT * FROM eleves WHERE moyenne >= 10");

  // Compter avant d'afficher
  $nb = mysqli_num_rows($res);
  echo "<p>Nombre d'admis : $nb</p>";

  if ($nb == 0) {
    echo "<p>Aucun élève admis pour le moment.</p>";
  } else {
    while ($e = mysqli_fetch_assoc($res)) {
      echo "<p>" . $e["nom"] . " — " . $e["moyenne"] . "</p>";
    }
  }
?>
\end{codephp}

\section{Insérer, modifier, supprimer depuis PHP}

Les requêtes \code{INSERT}, \code{UPDATE} et \code{DELETE} se lancent aussi avec
\code{mysqli\_query}. Le plus souvent, leurs valeurs proviennent d'un
\textbf{formulaire} (\code{\$\_POST} ou \code{\$\_GET}).

\subsection{Ajouter (INSERT) à partir d'un formulaire}

\begin{codehtml}
<!-- ajouter.php : le formulaire -->
<form method="post" action="ajouter.php">
  Nom :     <input type="text"   name="nom"><br>
  Classe :  <input type="text"   name="classe"><br>
  Moyenne : <input type="number" name="moyenne" step="0.01"><br>
  <input type="submit" value="Ajouter">
</form>
\end{codehtml}

\begin{codephp}
<?php
  require "connexion.php";

  // Le formulaire a-t-il été envoyé ?
  if (isset($_POST["nom"])) {
    // On récupère les données saisies
    $nom     = $_POST["nom"];
    $classe  = $_POST["classe"];
    $moyenne = $_POST["moyenne"];

    // On construit puis on exécute la requête INSERT
    $sql = "INSERT INTO eleves (nom, classe, moyenne)
            VALUES ('$nom', '$classe', $moyenne)";

    if (mysqli_query($cnx, $sql)) {
      echo "<p>Élève ajouté avec succès.</p>";
    } else {
      echo "<p>Erreur : " . mysqli_error($cnx) . "</p>";
    }
  }
?>
\end{codephp}

\begin{attention}
Dans la requête, le \code{\$nom} et la \code{\$classe} sont du texte : ils restent
entre apostrophes \code{'\$nom'}. La \code{\$moyenne} est un nombre : pas
d'apostrophes.
\end{attention}

\subsection{Modifier (UPDATE)}

\begin{codephp}
<?php
  require "connexion.php";

  $id      = $_POST["id"];       // id de l'élève à modifier
  $moyenne = $_POST["moyenne"];  // nouvelle moyenne

  $sql = "UPDATE eleves SET moyenne = $moyenne WHERE id = $id";

  if (mysqli_query($cnx, $sql)) {
    echo "<p>Moyenne mise à jour.</p>";
  } else {
    echo "<p>Erreur : " . mysqli_error($cnx) . "</p>";
  }
?>
\end{codephp}

\subsection{Supprimer (DELETE)}

On supprime une ligne précise grâce à sa clé primaire. L'\code{id} est souvent
passé dans l'URL (\code{supprimer.php?id=2}), donc lu dans \code{\$\_GET}.

\begin{codephp}
<?php
  require "connexion.php";

  $id = $_GET["id"];      // ex. supprimer.php?id=2

  $sql = "DELETE FROM eleves WHERE id = $id";

  if (mysqli_query($cnx, $sql)) {
    echo "<p>Élève supprimé.</p>";
  } else {
    echo "<p>Erreur : " . mysqli_error($cnx) . "</p>";
  }
?>
\end{codephp}

\section{Fermer la connexion : \texttt{mysqli\_close}}

À la fin du script, on libère la connexion au serveur.

\begin{codephp}
<?php
  mysqli_close($cnx);   // bonne pratique en fin de page
?>
\end{codephp}

\begin{remarque}
PHP ferme automatiquement la connexion à la fin de l'exécution de la page ;
\code{mysqli\_close} reste une bonne habitude, surtout pour les longs scripts.
\end{remarque}

\section{Application fil rouge : gestion des élèves}

Réunissons tout dans une \textbf{mini-application} de gestion des élèves : on
affiche la liste, on ajoute un élève, on en supprime un. Trois fichiers
suffisent.

\subsection{1. \texttt{connexion.php}}

\begin{codephp}
<?php
  $cnx = mysqli_connect("localhost", "root", "", "lycee");
  if (!$cnx) {
    die("Connexion impossible : " . mysqli_connect_error());
  }
  mysqli_set_charset($cnx, "utf8");
?>
\end{codephp}

\subsection{2. \texttt{liste.php} — afficher, ajouter, supprimer}

\begin{codephp}
<?php require "connexion.php"; ?>
<!DOCTYPE html>
<html lang="fr">
<head><meta charset="utf-8"><title>Gestion des élèves</title></head>
<body>
  <h1>Liste des élèves — Lycée de Maroua</h1>

<?php
  // --- AJOUT (formulaire envoyé en POST) ---
  if (isset($_POST["nom"])) {
    $nom     = $_POST["nom"];
    $classe  = $_POST["classe"];
    $moyenne = $_POST["moyenne"];
    $sql = "INSERT INTO eleves (nom, classe, moyenne)
            VALUES ('$nom', '$classe', $moyenne)";
    mysqli_query($cnx, $sql);
  }

  // --- SUPPRESSION (id passé dans l'URL) ---
  if (isset($_GET["suppr"])) {
    $id = $_GET["suppr"];
    mysqli_query($cnx, "DELETE FROM eleves WHERE id = $id");
  }
?>

  <!-- Formulaire d'ajout -->
  <form method="post" action="liste.php">
    <input type="text"   name="nom"     placeholder="Nom">
    <input type="text"   name="classe"  placeholder="Classe">
    <input type="number" name="moyenne" step="0.01" placeholder="Moyenne">
    <input type="submit" value="Ajouter">
  </form>

  <!-- Tableau des élèves -->
  <table border="1" cellpadding="6">
    <tr><th>Id</th><th>Nom</th><th>Classe</th><th>Moyenne</th><th>Action</th></tr>
<?php
  $res = mysqli_query($cnx, "SELECT * FROM eleves ORDER BY moyenne DESC");
  while ($e = mysqli_fetch_assoc($res)) {
    echo "<tr>";
    echo "<td>" . $e["id"]      . "</td>";
    echo "<td>" . $e["nom"]     . "</td>";
    echo "<td>" . $e["classe"]  . "</td>";
    echo "<td>" . $e["moyenne"] . "</td>";
    // Lien de suppression : on passe l'id dans l'URL
    echo "<td><a href='liste.php?suppr=" . $e["id"] . "'>Supprimer</a></td>";
    echo "</tr>";
  }
?>
  </table>

<?php mysqli_close($cnx); ?>
</body>
</html>
\end{codephp}
\fig{Une seule page gère l'affichage, l'ajout et la suppression : c'est une
mini-application CRUD complète.}

\begin{remarque}
\textbf{CRUD} résume les quatre opérations de base d'une application de gestion :
\emph{Create} (INSERT), \emph{Read} (SELECT), \emph{Update} (UPDATE) et
\emph{Delete} (DELETE). Tu sais désormais les réaliser toutes les quatre.
\end{remarque}

\begin{attention}
Ces exemples sont volontairement simples pour le programme TI. Dans une vraie
application, on \textbf{n'insère jamais} directement \code{\$\_POST} dans une
requête (risque d'\emph{injection SQL}) : on nettoie les données. Retiens
l'idée ; les techniques de sécurisation se voient en études supérieures.
\end{attention}

\section{Exercices}

\rubrique{Application}

\exo{Créer une table}\diff{1}
Écris la requête SQL qui crée une table \code{produits} avec les champs :
\code{id} (entier, clé primaire auto-incrémentée), \code{designation}
(texte, 50 caractères), \code{prix} (nombre décimal), \code{stock} (entier).

\exo{Insérer des lignes}\diff{1}
Insère dans la table \code{produits} deux articles d'une boutique de Yaoundé :
un « Cahier 200 pages » à 350 FCFA (stock 120) et un « Stylo bleu Bic »
à 100 FCFA (stock 500).

\exo{Sélectionner les admis}\diff{1}
Écris la requête qui affiche le nom et la moyenne de tous les élèves
\textbf{admis} (moyenne supérieure ou égale à 10), triés de la meilleure
moyenne à la moins bonne.

\exo{Mettre à jour}\diff{1}
Écris la requête qui augmente de 50 le stock du produit dont l'\code{id} vaut 1.

\exo{Supprimer}\diff{1}
Écris la requête qui supprime tous les produits en rupture de stock
(\code{stock = 0}).

\exo{Afficher une table en HTML}\diff{2}
Écris le script PHP qui se connecte à la base \code{lycee}, lit la table
\code{produits} et affiche \code{designation}, \code{prix} et \code{stock}
dans un tableau HTML.

\exo{Compter les enregistrements}\diff{2}
Écris le script PHP qui affiche le \textbf{nombre} d'élèves de la classe
\code{'Tle TI'} contenus dans la table \code{eleves}.

\rubrique{Entraînement}

\exo{Recherche par nom}\diff{2}
Un formulaire envoie en POST un champ \code{recherche}. Écris le script PHP qui
affiche tous les élèves dont le nom \textbf{contient} le texte saisi (on
utilisera \code{LIKE} et le caractère \code{\%}).

\exo{Moyenne de la classe}\diff{2}
Écris une requête SQL qui calcule la \textbf{moyenne générale} (moyenne des
moyennes) de tous les élèves de la table \code{eleves} (fonction \code{AVG}).

\exo{Formulaire de modification}\diff{3}
Écris un script \code{modifier.php} qui reçoit en POST un \code{id} et une
nouvelle \code{moyenne}, met à jour l'élève correspondant, puis affiche un
message de confirmation.

\section{Corrigés}

\corrige{de l'exercice 1}
\begin{codesql}
CREATE TABLE produits (
  id          INT AUTO_INCREMENT PRIMARY KEY,
  designation VARCHAR(50) NOT NULL,
  prix        DECIMAL(8,2) DEFAULT 0,
  stock       INT DEFAULT 0
);
\end{codesql}

\corrige{de l'exercice 2}
\begin{codesql}
INSERT INTO produits (designation, prix, stock)
VALUES ('Cahier 200 pages', 350, 120);

INSERT INTO produits (designation, prix, stock)
VALUES ('Stylo bleu Bic', 100, 500);
\end{codesql}

\corrige{de l'exercice 3}
\begin{codesql}
SELECT nom, moyenne FROM eleves
WHERE moyenne >= 10
ORDER BY moyenne DESC;
\end{codesql}

\corrige{de l'exercice 4}
\begin{codesql}
UPDATE produits SET stock = stock + 50 WHERE id = 1;
\end{codesql}

\corrige{de l'exercice 5}
\begin{codesql}
DELETE FROM produits WHERE stock = 0;
\end{codesql}

\corrige{de l'exercice 6}
\begin{codephp}
<?php
  $cnx = mysqli_connect("localhost", "root", "", "lycee");
  if (!$cnx) {
    die("Connexion impossible : " . mysqli_connect_error());
  }
  $res = mysqli_query($cnx, "SELECT * FROM produits");
?>
<table border="1">
  <tr><th>Désignation</th><th>Prix (FCFA)</th><th>Stock</th></tr>
<?php
  while ($p = mysqli_fetch_assoc($res)) {
    echo "<tr>";
    echo "<td>" . $p["designation"] . "</td>";
    echo "<td>" . $p["prix"]        . "</td>";
    echo "<td>" . $p["stock"]       . "</td>";
    echo "</tr>";
  }
  mysqli_close($cnx);
?>
</table>
\end{codephp}

\corrige{de l'exercice 7}
\begin{codephp}
<?php
  $cnx = mysqli_connect("localhost", "root", "", "lycee");
  $res = mysqli_query($cnx, "SELECT * FROM eleves WHERE classe = 'Tle TI'");
  $nb  = mysqli_num_rows($res);
  echo "<p>La classe de Tle TI compte $nb élève(s).</p>";
  mysqli_close($cnx);
?>
\end{codephp}

\corrige{de l'exercice 8}
On colle le texte saisi entre deux \code{\%} : \code{LIKE '\%mot\%'} cherche le
mot n'importe où dans le nom.
\begin{codephp}
<?php
  $cnx = mysqli_connect("localhost", "root", "", "lycee");
  $mot = $_POST["recherche"];

  $sql = "SELECT * FROM eleves WHERE nom LIKE '%$mot%'";
  $res = mysqli_query($cnx, $sql);

  if (mysqli_num_rows($res) == 0) {
    echo "<p>Aucun élève trouvé.</p>";
  } else {
    while ($e = mysqli_fetch_assoc($res)) {
      echo "<p>" . $e["nom"] . " (" . $e["classe"] . ")</p>";
    }
  }
  mysqli_close($cnx);
?>
\end{codephp}

\corrige{de l'exercice 9}
\begin{codesql}
SELECT AVG(moyenne) AS moyenne_generale FROM eleves;
\end{codesql}

\corrige{de l'exercice 10}
\begin{codephp}
<?php
  // modifier.php
  $cnx = mysqli_connect("localhost", "root", "", "lycee");
  if (!$cnx) {
    die("Connexion impossible : " . mysqli_connect_error());
  }

  if (isset($_POST["id"])) {
    $id      = $_POST["id"];
    $moyenne = $_POST["moyenne"];

    $sql = "UPDATE eleves SET moyenne = $moyenne WHERE id = $id";

    if (mysqli_query($cnx, $sql)) {
      echo "<p>Élève n°$id mis à jour : nouvelle moyenne = $moyenne.</p>";
    } else {
      echo "<p>Erreur : " . mysqli_error($cnx) . "</p>";
    }
  }
  mysqli_close($cnx);
?>
\end{codephp}

\begin{aretenir}
Pour interfacer une base MySQL en PHP procédural : \textbf{(1)} se connecter avec
\code{mysqli\_connect} (et \code{die} + \code{mysqli\_connect\_error} en cas
d'échec) ; \textbf{(2)} exécuter la requête avec \code{mysqli\_query} ;
\textbf{(3)} pour un \code{SELECT}, parcourir le résultat avec
\code{while (\$l = mysqli\_fetch\_assoc(\$res))} ; \textbf{(4)} fermer avec
\code{mysqli\_close}. En SQL, n'oublie \textbf{jamais} le \code{WHERE} dans un
\code{UPDATE} ou un \code{DELETE}.
\end{aretenir}
